% Section 3: PPO
\section{Proximal Policy Optimization (PPO)}
\label{sec:ppo}

\subsection{Algorithm Overview and Motivation}

PPO~\cite{schulman2017ppo}, proposed by Schulman et al.\ in 2017, remains the most widely adopted RL algorithm for LLM alignment. It was the primary algorithm behind InstructGPT~\cite{ouyang2022instructgpt} and early versions of ChatGPT. PPO builds upon the policy gradient framework with a clipped surrogate objective that constrains policy updates to a trust region, preventing destructively large updates while maintaining sample efficiency.

The key insight behind PPO is that vanilla policy gradient methods suffer from high variance and unstable updates. Trust Region Policy Optimization (TRPO) addressed this with a hard KL constraint, but its second-order optimization is computationally expensive. PPO achieves similar stability with first-order optimization by clipping the probability ratio.

\subsection{Clipped Surrogate Objective}

The core PPO objective is the clipped surrogate function. Let $r_t(\theta) = \frac{\pi_\theta(a_t|s_t)}{\pi_{\theta_{\text{old}}}(a_t|s_t)}$ denote the probability ratio between the current and old policies. The clipped objective is:
\begin{equation}
L^{\text{CLIP}}(\theta) = \E_t\!\left[\min\!\left(r_t(\theta)\hat{A}_t,\ \text{clip}(r_t(\theta), 1{-}\varepsilon, 1{+}\varepsilon)\hat{A}_t\right)\right]
\label{eq:ppo_clip}
\end{equation}
where $\varepsilon$ is a hyperparameter (typically 0.2) controlling the clipping range.

\textbf{Interpretation:} When the advantage $\hat{A}_t > 0$ (the action was better than expected), the objective encourages increasing $r_t$ but clips it at $1 + \varepsilon$, preventing excessive policy updates. When $\hat{A}_t < 0$, it discourages the action but clips $r_t$ at $1 - \varepsilon$. This creates a ``flat'' region in the objective that provides robustness to large, potentially catastrophic updates.

\subsection{Value Function and Entropy Bonus}

PPO maintains a separate value network $V_\phi(s)$ to estimate expected returns. The value function loss uses a mean-squared error objective:
\begin{equation}
L^{VF}(\phi) = \E_t\left[\left(V_\phi(s_t) - \hat{V}_t^{\text{targ}}\right)^2\right]
\end{equation}
where $\hat{V}_t^{\text{targ}}$ is the target value computed from discounted returns.

To encourage exploration and prevent premature convergence to deterministic policies, PPO adds an entropy bonus:
\begin{equation}
S[\pi_\theta](s_t) = -\sum_{a} \pi_\theta(a|s_t)\log\pi_\theta(a|s_t)
\end{equation}

\subsection{Combined PPO Objective}

The full PPO objective combines all three components:
\begin{equation}
\boxed{L^{\text{PPO}}(\theta, \phi) = \E_t\left[L^{\text{CLIP}}(\theta) - c_1 L^{VF}(\phi) + c_2 S[\pi_\theta]\right]}
\label{eq:ppo}
\end{equation}
where $c_1 = 0.5$ and $c_2 = 0.01$ are typical coefficient values. The expectation is taken over a batch of trajectories collected by the current policy.

\subsection{Generalized Advantage Estimation (GAE)}

The advantage function $\hat{A}_t$ is estimated using GAE~\cite{schulman2016gae}, which provides a bias-variance trade-off controlled by the parameter $\lambda$:
\begin{equation}
\hat{A}_t^{\text{GAE}} = \sum_{l=0}^{\infty}(\gamma\lambda)^l\delta_{t+l}^{V}
\end{equation}
where $\delta_t^{V} = r_t + \gamma V_\phi(s_{t+1}) - V_\phi(s_t)$ is the temporal difference (TD) error. When $\lambda = 0$, GAE reduces to the one-step TD error (low variance, high bias); when $\lambda = 1$, it reduces to Monte Carlo returns (low bias, high variance).

\subsection{Architectural Requirements}

PPO requires \textbf{four model components} to be simultaneously loaded into GPU memory:
\begin{enumerate}[leftmargin=*, itemsep=1pt]
\item \textbf{Actor} (policy $\pi_\theta$): The LLM being trained.
\item \textbf{Critic} (value $V_\phi$): A separate network of comparable size to the Actor.
\item \textbf{Reference Model} ($\pi_{\text{ref}}$): Frozen copy for KL computation (optional but common in RLHF).
\item \textbf{Reward Model} ($r_\psi$): The trained reward function.
\end{enumerate}

This makes PPO the most memory-intensive algorithm, requiring approximately $4\times$ the memory of a single LLM forward pass. For a 7B parameter model, this typically requires 4--8 A100 GPUs with model parallelism.

\subsection{Strengths and Limitations}

\textbf{Strengths:} Well-understood convergence properties; robust to hyperparameter choices; compatible with any reward function; strong theoretical foundation from policy gradient theory.

\textbf{Limitations:} High memory overhead from the value network; training instability when the reward signal is sparse or noisy; requires careful reward model training; the value network adds significant engineering complexity.
